\documentclass[10pt]{extarticle}
\usepackage[a4paper,margin=1.5cm]{geometry}
\usepackage[T1]{fontenc}
\usepackage{amsmath,amssymb,mathtools}
\usepackage{graphicx}
\usepackage[export]{adjustbox}
\usepackage{multicol}
\usepackage{enumitem}
\usepackage{booktabs,array,tabularx}
\setlength{\columnsep}{0.7cm}
\setlength{\parindent}{0pt}
\setlength{\parskip}{3pt}
\pagestyle{empty}
\makeatletter
\newcommand{\qfigdef}[4]{\@namedef{qf@file@#1}{#2}\@namedef{qf@w@#1}{#3}\@namedef{qf@wide@#1}{#4}}
\newcommand{\qf@check}[1]{\@ifundefined{qf@file@#1}{\PackageError{score2}{Unknown figure id #1}{}}{}}
\newcommand{\qf@img}[4][]{\includegraphics[width=#2\dimexpr\@nameuse{qf@w@#4}pt\relax,max width=#3,max height=0.45\textheight,#1]{\@nameuse{qf@file@#4}}}
\newif\ifqf@full
\newif\ifqf@boxed
\newsavebox{\qf@box}
\newcommand{\qf@open}{\global\qf@fulltrue\par\global\qf@boxedtrue
  \setbox\qf@box\vbox\bgroup\hsize\textwidth\linewidth\textwidth}
\newcommand{\qf@emit}[1]{\par\egroup\global\qf@boxedfalse
  \ifdim\dimexpr\ht\qf@box+\dp\qf@box\relax<#1\textheight\noindent\usebox{\qf@box}\par\else\unvbox\qf@box\fi}
\newcommand{\qfullsplit}{\ifqf@boxed\qf@emit{0.85}\fi}
\newcommand{\qf@close}{\ifqf@boxed\qf@emit{0.45}\fi\global\qf@fullfalse}
\newcommand{\qfullbegin}{\end{multicols}\qf@open}
\newcommand{\qfullend}{\qf@close\begin{multicols}{2}}
\newcommand{\qfullstart}{\qf@open}
\newcommand{\qfulldone}{\qf@close}
\newcommand{\qfig}[2][1]{\qf@check{#2}%
  \ifqf@full\qf@blk{#1}{#2}\else
  \ifnum\@nameuse{qf@wide@#2}=1 \expandafter\@firstoftwo\else\expandafter\@secondoftwo\fi
  {\end{multicols}\begin{center}\qf@img{#1}{\textwidth}{#2}\end{center}\begin{multicols}{2}}%
  {\qf@blk{#1}{#2}}\fi}
\newcommand{\qf@blk}[2]{\par\smallskip{\centering\qf@img{#1}{\linewidth}{#2}\par}\smallskip}
\newcommand{\qfigi}[2][1]{\qf@check{#2}\qf@img[valign=c]{#1}{\linewidth}{#2}}
\makeatother
\newcommand{\papersec}[1]{\par\bigskip{\centering\Large\bfseries PAPER-#1\par}\medskip}
\newcommand{\partsec}[2]{\par\medskip\noindent{\large\bfseries PART-#1: #2}\par\smallskip}
\newcommand{\sect}[1]{\par\smallskip\noindent\textbf{#1}\par}
\newcommand{\qn}[1]{\par\smallskip\noindent\textbf{#1.}~}
\newcommand{\paraq}[2]{\par\smallskip\noindent\textbf{Paragraph for Questions #1 and #2}\par}
\newenvironment{qoptions}{\begin{itemize}[leftmargin=1.8em,itemsep=1pt,topsep=2pt,label={}]}{\end{itemize}}
\newcommand{\opt}[2]{\item[(#1)] #2}
\newcolumntype{L}[1]{>{\raggedright\arraybackslash}p{#1}}
\newcommand{\fwstart}{\end{multicols}\par\noindent\begin{minipage}{\textwidth}\small\setlength{\tabcolsep}{4pt}\renewcommand{\arraystretch}{1.5}}
\newcommand{\fwend}{\end{minipage}\par\medskip\begin{multicols}{2}}
% figure registry (generated by score2_tools.py finalize)
\qfigdef{p1_phy_s1_q04_1}{Pre-Score (2023-24) - Test 08 (03 Dec 2023)-img/p1_phy_s1_q04_1.png}{169.5}{0}
\qfigdef{p1_phy_s1_q04_2}{Pre-Score (2023-24) - Test 08 (03 Dec 2023)-img/p1_phy_s1_q04_2.png}{169.0}{0}
\qfigdef{p1_phy_s2_q05_1}{Pre-Score (2023-24) - Test 08 (03 Dec 2023)-img/p1_phy_s2_q05_1.png}{83.2}{0}
\qfigdef{p1_phy_s2_q06_1}{Pre-Score (2023-24) - Test 08 (03 Dec 2023)-img/p1_phy_s2_q06_1.png}{115.8}{0}
\qfigdef{p1_phy_s2_q07_1}{Pre-Score (2023-24) - Test 08 (03 Dec 2023)-img/p1_phy_s2_q07_1.png}{160.9}{0}
\qfigdef{p1_phy_s3_q08_1}{Pre-Score (2023-24) - Test 08 (03 Dec 2023)-img/p1_phy_s3_q08_1.png}{161.3}{0}
\qfigdef{p1_phy_s3_q09_1}{Pre-Score (2023-24) - Test 08 (03 Dec 2023)-img/p1_phy_s3_q09_1.png}{107.0}{0}
\qfigdef{p1_phy_s3_q09_2}{Pre-Score (2023-24) - Test 08 (03 Dec 2023)-img/p1_phy_s3_q09_2.png}{96.0}{0}
\qfigdef{p1_phy_s3_q09_3}{Pre-Score (2023-24) - Test 08 (03 Dec 2023)-img/p1_phy_s3_q09_3.png}{113.4}{0}
\qfigdef{p1_phy_s3_q09_4}{Pre-Score (2023-24) - Test 08 (03 Dec 2023)-img/p1_phy_s3_q09_4.png}{113.6}{0}
\qfigdef{p1_phy_s3_q10_1}{Pre-Score (2023-24) - Test 08 (03 Dec 2023)-img/p1_phy_s3_q10_1.png}{124.2}{0}
\qfigdef{p1_phy_s4_q01_1}{Pre-Score (2023-24) - Test 08 (03 Dec 2023)-img/p1_phy_s4_q01_1.png}{162.3}{0}
\qfigdef{p1_phy_s4_q02_1}{Pre-Score (2023-24) - Test 08 (03 Dec 2023)-img/p1_phy_s4_q02_1.png}{105.3}{0}
\qfigdef{p1_phy_s4_q05_1}{Pre-Score (2023-24) - Test 08 (03 Dec 2023)-img/p1_phy_s4_q05_1.png}{161.1}{0}
\qfigdef{p1_chem_s1_q03_A}{Pre-Score (2023-24) - Test 08 (03 Dec 2023)-img/p1_chem_s1_q03_A.png}{108.9}{0}
\qfigdef{p1_chem_s1_q03_B}{Pre-Score (2023-24) - Test 08 (03 Dec 2023)-img/p1_chem_s1_q03_B.png}{242.9}{0}
\qfigdef{p1_chem_s1_q03_C}{Pre-Score (2023-24) - Test 08 (03 Dec 2023)-img/p1_chem_s1_q03_C.png}{245.2}{0}
\qfigdef{p1_chem_s2_q07_1}{Pre-Score (2023-24) - Test 08 (03 Dec 2023)-img/p1_chem_s2_q07_1.png}{240.0}{0}
\qfigdef{p1_chem_s3_q09_1}{Pre-Score (2023-24) - Test 08 (03 Dec 2023)-img/p1_chem_s3_q09_1.png}{45.1}{0}
\qfigdef{p1_chem_s3_q09_2}{Pre-Score (2023-24) - Test 08 (03 Dec 2023)-img/p1_chem_s3_q09_2.png}{22.3}{0}
\qfigdef{p1_chem_s3_q09_3}{Pre-Score (2023-24) - Test 08 (03 Dec 2023)-img/p1_chem_s3_q09_3.png}{52.4}{0}
\qfigdef{p1_chem_s3_q09_4}{Pre-Score (2023-24) - Test 08 (03 Dec 2023)-img/p1_chem_s3_q09_4.png}{29.7}{0}
\qfigdef{p1_chem_s3_q10_1}{Pre-Score (2023-24) - Test 08 (03 Dec 2023)-img/p1_chem_s3_q10_1.png}{77.9}{0}
\qfigdef{p1_chem_s3_q10_2}{Pre-Score (2023-24) - Test 08 (03 Dec 2023)-img/p1_chem_s3_q10_2.png}{112.3}{0}
\qfigdef{p1_chem_s3_q10_3}{Pre-Score (2023-24) - Test 08 (03 Dec 2023)-img/p1_chem_s3_q10_3.png}{108.0}{0}
\qfigdef{p1_chem_s3_q10_4}{Pre-Score (2023-24) - Test 08 (03 Dec 2023)-img/p1_chem_s3_q10_4.png}{94.6}{0}
\qfigdef{p1_chem_s4_q06_1}{Pre-Score (2023-24) - Test 08 (03 Dec 2023)-img/p1_chem_s4_q06_1.png}{200.6}{0}
\qfigdef{p2_phy_s1_q01_1}{Pre-Score (2023-24) - Test 08 (03 Dec 2023)-img/p2_phy_s1_q01_1.png}{108.3}{0}
\qfigdef{p2_phy_s1_q03_1}{Pre-Score (2023-24) - Test 08 (03 Dec 2023)-img/p2_phy_s1_q03_1.png}{134.7}{0}
\qfigdef{p2_phy_s1_q04_1}{Pre-Score (2023-24) - Test 08 (03 Dec 2023)-img/p2_phy_s1_q04_1.png}{145.4}{0}
\qfigdef{p2_phy_s2_q05_1}{Pre-Score (2023-24) - Test 08 (03 Dec 2023)-img/p2_phy_s2_q05_1.png}{76.9}{0}
\qfigdef{p2_phy_s2_q07_1}{Pre-Score (2023-24) - Test 08 (03 Dec 2023)-img/p2_phy_s2_q07_1.png}{195.1}{0}
\qfigdef{p2_phy_s3_para1_2_1}{Pre-Score (2023-24) - Test 08 (03 Dec 2023)-img/p2_phy_s3_para1_2_1.png}{102.7}{0}
\qfigdef{p2_phy_s3_para3_4_1}{Pre-Score (2023-24) - Test 08 (03 Dec 2023)-img/p2_phy_s3_para3_4_1.png}{123.2}{0}
\qfigdef{p2_phy_s4_q06_1}{Pre-Score (2023-24) - Test 08 (03 Dec 2023)-img/p2_phy_s4_q06_1.png}{142.0}{0}
\qfigdef{p2_phy_s4_q07_1}{Pre-Score (2023-24) - Test 08 (03 Dec 2023)-img/p2_phy_s4_q07_1.png}{115.4}{0}
\qfigdef{p2_phy_s4_q09_1}{Pre-Score (2023-24) - Test 08 (03 Dec 2023)-img/p2_phy_s4_q09_1.png}{85.0}{0}
\qfigdef{p2_phy_s4_q10_1}{Pre-Score (2023-24) - Test 08 (03 Dec 2023)-img/p2_phy_s4_q10_1.png}{119.3}{0}
\qfigdef{p2_chem_s1_q01_1}{Pre-Score (2023-24) - Test 08 (03 Dec 2023)-img/p2_chem_s1_q01_1.png}{180.8}{0}
\qfigdef{p2_chem_s1_q01_A}{Pre-Score (2023-24) - Test 08 (03 Dec 2023)-img/p2_chem_s1_q01_A.png}{75.1}{0}
\qfigdef{p2_chem_s1_q01_B}{Pre-Score (2023-24) - Test 08 (03 Dec 2023)-img/p2_chem_s1_q01_B.png}{57.5}{0}
\qfigdef{p2_chem_s1_q01_C}{Pre-Score (2023-24) - Test 08 (03 Dec 2023)-img/p2_chem_s1_q01_C.png}{55.4}{0}
\qfigdef{p2_chem_s1_q01_D}{Pre-Score (2023-24) - Test 08 (03 Dec 2023)-img/p2_chem_s1_q01_D.png}{55.6}{0}
\qfigdef{p2_chem_s1_q03_1}{Pre-Score (2023-24) - Test 08 (03 Dec 2023)-img/p2_chem_s1_q03_1.png}{280.6}{1}
\qfigdef{p2_chem_s2_q05_A}{Pre-Score (2023-24) - Test 08 (03 Dec 2023)-img/p2_chem_s2_q05_A.png}{82.8}{0}
\qfigdef{p2_chem_s2_q05_B}{Pre-Score (2023-24) - Test 08 (03 Dec 2023)-img/p2_chem_s2_q05_B.png}{105.5}{0}
\qfigdef{p2_chem_s2_q05_C}{Pre-Score (2023-24) - Test 08 (03 Dec 2023)-img/p2_chem_s2_q05_C.png}{76.1}{0}
\qfigdef{p2_chem_s2_q05_D}{Pre-Score (2023-24) - Test 08 (03 Dec 2023)-img/p2_chem_s2_q05_D.png}{113.2}{0}
\qfigdef{p2_chem_s2_q06_1}{Pre-Score (2023-24) - Test 08 (03 Dec 2023)-img/p2_chem_s2_q06_1.png}{134.4}{0}
\qfigdef{p2_chem_s3_para1_2_1}{Pre-Score (2023-24) - Test 08 (03 Dec 2023)-img/p2_chem_s3_para1_2_1.png}{109.0}{0}
\qfigdef{p2_chem_s3_q02_1}{Pre-Score (2023-24) - Test 08 (03 Dec 2023)-img/p2_chem_s3_q02_1.png}{131.8}{0}
\qfigdef{p2_chem_s4_q06_1}{Pre-Score (2023-24) - Test 08 (03 Dec 2023)-img/p2_chem_s4_q06_1.png}{83.7}{0}
\qfigdef{p2_chem_s4_q07_1}{Pre-Score (2023-24) - Test 08 (03 Dec 2023)-img/p2_chem_s4_q07_1.png}{174.1}{0}
\qfigdef{p2_chem_s4_q10_1}{Pre-Score (2023-24) - Test 08 (03 Dec 2023)-img/p2_chem_s4_q10_1.png}{156.2}{0}
\raggedcolumns\makeatletter\@beginparpenalty=10000\makeatother

\begin{document}
\begin{center}{\Large\bfseries Pre-Score (2023-24) --- Test 08 (03 Dec 2023)}\end{center}
\begin{multicols}{2}
\papersec{1}
\partsec{1}{PHYSICS}
\sect{SECTION-I (i)}
\qn{1} A plane electromagnetic wave of wavelength $\lambda$ has an intensity $I$. It is propagating along the positive $Y$-direction. The allowed expressions for the electric and magnetic fields are given by
\begin{qoptions}
\opt{A}{$\vec E=\sqrt{\dfrac{2I}{\epsilon_0c}}\cos\!\left[\dfrac{2\pi}{\lambda}(y-ct)\right]\hat k;\ \vec B=+\dfrac1cE\,\hat i$}
\opt{B}{$\vec E=\sqrt{\dfrac{I}{\epsilon_0c}}\cos\!\left[\dfrac{2\pi}{\lambda}(y-ct)\right]\hat k;\ \vec B=+\dfrac1cE\,\hat i$}
\opt{C}{$\vec E=\sqrt{\dfrac{2I}{\epsilon_0c}}\cos\!\left[\dfrac{2\pi}{\lambda}(y-ct)\right]\hat k;\ \vec B=-\dfrac1cE\,\hat i$}
\opt{D}{$\vec E=\sqrt{\dfrac{I}{\epsilon_0c}}\cos\!\left[\dfrac{2\pi}{\lambda}(y-ct)\right]\hat i;\ \vec B=\dfrac1cE\,\hat k$}
\end{qoptions}
\qn{2} Assume that sun rises at 6 : 00 AM \& sets at 6 : 00 PM. It moves such that length of shadow of a vertical object at 12 : 00 Noon is zero. An observer observes the reflected rays from surface of water $\left(\mu=\dfrac43\right)$ using a polaroid. At what approximate time will he see no light at a particular orientation of polaroid?
\begin{qoptions}
\opt{A}{10 : 00 AM}
\opt{B}{11 : 00 AM}
\opt{C}{9 : 30 AM}
\opt{D}{8 : 30 AM}
\end{qoptions}
\qn{3} Longitudinal Standing wave produced in a metal rod of length 1m fixed at the left end is represented by the equation\\
$y=10^{-6}\sin\dfrac{\pi x}{2}\sin200\pi t$\\
where $x$ is in metre and $t$ is in seconds.\\
The maximum tensile stress at the midpoint of the rod is : (Young's modulus of material of rod $=10^{12}$ N/m$^2$)
\begin{qoptions}
\opt{A}{$\dfrac{\pi}{\sqrt2}\times10^6$N/m$^2$}
\opt{B}{Zero}
\opt{C}{$\dfrac{\pi}{2\sqrt2}\times10^6$N/m$^2$}
\opt{D}{$\dfrac\pi2\times10^6$N/m$^2$}
\end{qoptions}
\qn{4} In a specially designed vernier calipers, 10 divisions of the vernier scale are equal to 4 divisions of the main scale. In the following figure, the zero error and a measurement are shown. In the zero error, the 5$^{th}$ and the 10$^{th}$ divisions of the vernier coincides with 1 mm and 3 mm mark on the main scale respectively and in the measurement, the 2$^{nd}$ and the 7$^{th}$ divisions of the vernier scale coincide with the 23 mm and 25 mm mark on the main scale respectively as shown in the figure. Report the length measured with due regards to the error.
\qfig{p1_phy_s1_q04_1}
\qfig{p1_phy_s1_q04_2}
\begin{qoptions}
\opt{A}{$(21.2\pm0.2)$ mm}
\opt{B}{$(23.2\pm0.2)$ mm}
\opt{C}{$(23.8\pm0.4)$ mm}
\opt{D}{$(23.4\pm0.6)$ mm}
\end{qoptions}
\sect{SECTION-I (ii)}
\qn{5} An uncharged conducting sphere of radius $R$, has a hollow cavity sphere of radius $\dfrac R2$, as shown. Which of the following is/are correct about electrostatic energy. A charge $Q$ is placed at the center of the cavity.
\qfig{p1_phy_s2_q05_1}
\begin{qoptions}
\opt{A}{Interaction energy of induced charges with charge placed at centre of cavity is $-\left(\frac{1}{4\pi\varepsilon_0}\right)\left(\frac{Q^2}{R}\right)$}
\opt{B}{Interaction energy of induced charges with charge placed at centre of cavity is $-\frac{1}{2\pi\varepsilon_0}\frac{Q^2}{R}$}
\opt{C}{Electrostatic energy outside sphere is $\frac{1}{4\pi\varepsilon_0}\left(\frac{Q^2}{2R}\right)$}
\opt{D}{Electrostatic energy outside sphere is $\left(\frac{1}{4\pi\varepsilon_0}\right)\left(\frac{Q^2}{R}\right)$}
\end{qoptions}
\qn{6} An infinite wire, placed along $z$-axis has current $I_1$ in positive $z$-direction. A conducting rod PQ placed in $xy$ plane parallel to $y$-axis has current $I_2$ in positive $y$-direction. The ends of the rod subtend angles $37^\circ$ and $53^\circ$ at the origin with positive $x$-axis as shown in figure. The rod is at a distance $\ell$ from the origin. Then choose the correct option(s).
\qfig{p1_phy_s2_q06_1}
\begin{qoptions}
\opt{A}{The direction of net force acting on the wire PQ is along the negative $z$-axis.}
\opt{B}{The direction of net force acting on the wire PQ is along the positive $z$-axis.}
\opt{C}{The magnitude of net force acting on the rod PQ is $\dfrac{\mu_0I_1I_2}{4\pi}\ell n\left(\dfrac{16}{9}\right)$}
\opt{D}{The magnitude of net force acting on the rod PQ is $\dfrac{\mu_0I_1I_2}{2\pi}\ell n\left(\dfrac{16}{9}\right)$}
\end{qoptions}
\qn{7} Two discs A and B each of mass 2 kg and radius 0.5 m are rigidly fixed to the ends of a thin shaft S, which is supported in horizontal position with the help of two smooth bearings at points $P_1$ and $P_2$. Two blocks of masses 2 kg and 1 kg are connected to the light cords wrapped around the discs as shown in the figure. There is no slipping between the cord and the discs. If the system is released from rest from the position shown, Then mark the correct statement(s). (Take : $g=10$ m/s$^2$)
\qfig{p1_phy_s2_q07_1}
\begin{qoptions}
\opt{A}{Power developed by tension in cord connected to 2 kg block is negative.}
\opt{B}{Acceleration of each block is 2 m/s$^2$.}
\opt{C}{Total work done by the tension on the discs is zero during the motion.}
\opt{D}{Acceleration of each block is 1 m/s$^2$.}
\end{qoptions}
\sect{SECTION-I (iii)}
\qn{8} The below diagram shows the set up for a variation of a Young's Double Slit Experiment. Plane parallel beams of two wavelengths $\lambda_1=2000$ \AA{} and $\lambda_2=7500$ \AA{} are incident on two slits $S_1$ and $S_2$ (assume $S_1$ and $S_2$ to be of identical dimension) the superposed waves from $S_1$ and $S_2$ are then incident on a second opaque screen with two slits $S_3$ and $S_4$ as shown. Superposed waves from $S_3$ and $S_4$ are then incident on a far-away screen $S$.\\
Take $\dfrac Dd=2\times10^3$, $y=0.5$ mm and assume all slits to be of identical size and both beams $\lambda_1=2000$ \AA{} and $\lambda_2=7500$ \AA{} of identical intensity $I_0$.\\
(i.e. from a single slit both $S_1$ and $S_2$ produce an intensity $I_0$ on $S_3S_4$ plane when considered one at a time).\\
(Also assume that $S_3$ and $S_4$ produce the same intensity on the screen $S$ considered one at a time as they receive from $S_1-S_2$ combination).
\qfig{p1_phy_s3_q08_1}
\par\nobreak
{\small\setlength{\tabcolsep}{3pt}\renewcommand{\arraystretch}{1.4}
\begin{tabular}{|L{0.8cm}|L{3.8cm}|L{0.8cm}|L{1.5cm}|}\hline
\multicolumn{2}{|c|}{\textbf{List-I}} & \multicolumn{2}{c|}{\textbf{List-II}}\\\hline
(P) & Intensity of light of wavelength $\lambda_1=2000$\AA{} at slit $S_3$ & (1) & $I_0$\\\hline
(Q) & Intensity of light of wavelength $\lambda_2=7500$\AA{} at slit $S_3$ & (2) & $2I_0$\\\hline
(R) & The maximum intensity for $\lambda_1=2000$\AA{} received at screen S & (3) & $9I_0$\\\hline
(S) & The maximum intensity for $\lambda_2=7500$\AA{} received at screen S & (4) & $\approx11.7\,I_0$\\\hline
 & & (5) & $4I_0$\\\hline
\end{tabular}}
\begin{qoptions}
\opt{A}{P $\to$ 2; Q $\to$ 1; R $\to$ 4; S $\to$ 3}
\opt{B}{P $\to$ 1; Q $\to$ 2; R $\to$ 4; S $\to$ 3}
\opt{C}{P $\to$ 2; Q $\to$ 4; R $\to$ 1; S $\to$ 3}
\opt{D}{P $\to$ 2; Q $\to$ 1; R $\to$ 3; S $\to$ 4}
\end{qoptions}
\fwstart
\qn{9} Two moles of an ideal monoatomic gas is taken through four different cyclic processes as depicted in the List-I and List-II gives the total work done by the gas during each cyclic process. (Take : $\ell$n2 $=0.7$)\par\nobreak
\begin{center}\begin{tabular}{|L{0.8cm}|L{5.2cm}|L{0.8cm}|L{1.8cm}|}\hline
\multicolumn{2}{|c|}{\textbf{List-I}} & \multicolumn{2}{c|}{\textbf{List-II}}\\\hline
(P) & In the process AB, $PT^{-1/2}=$ constant\newline \qfigi{p1_phy_s3_q09_1} & (1) & 480R\\\hline
(Q) & In the process BC, $PT^{-2}=$ constant\newline \qfigi{p1_phy_s3_q09_2} & (2) & 240R\\\hline
(R) & In the process BC, $VT^{-2}=$ constant\newline \qfigi{p1_phy_s3_q09_3} & (3) & 600R\\\hline
(S) & In the process CA, $VT^{-2}=$ constant\newline \qfigi{p1_phy_s3_q09_4} & (4) & 540R\\\hline
 & & (5) & 300R\\\hline
\end{tabular}\end{center}
Which one of the following options is correct ?
\begin{qoptions}
\opt{A}{P $\to$ 5; Q $\to$ 1; R $\to$ 3; S $\to$ 4}
\opt{B}{P $\to$ 5; Q $\to$ 1; R $\to$ 4; S $\to$ 2}
\opt{C}{P $\to$ 1; Q $\to$ 3; R $\to$ 4; S $\to$ 2}
\opt{D}{P $\to$ 2; Q $\to$ 4; R $\to$ 5; S $\to$ 1}
\end{qoptions}
\fwend
\qn{10} A conducting sphere of radius $a$ is placed inside a thin conducting shell of radius $4a$. Now a dielectric shell of inner radius $2a$ and outer radius $3a$ is placed between sphere and shell concentrically. A is the outer surface of the sphere, B is the inner surface of the dielectric shell, C is the outer surface of dielectric shell and D is the inner surface of conducting shell. The charge $+q$ is given to the outer surface of sphere A.
\qfig{p1_phy_s3_q10_1}
\par\nobreak
{\small\setlength{\tabcolsep}{3pt}\renewcommand{\arraystretch}{1.5}
\begin{tabular}{|L{1.6cm}|L{0.8cm}|L{3.4cm}|}\hline
\multicolumn{1}{|c|}{\textbf{List-I}} & \multicolumn{2}{c|}{\textbf{List-II}}\\\hline
(P) Surface A & (1) & Charge is $-q$\\\hline
(Q) Surface B & (2) & Charge is $+q$\\\hline
(R) Surface C & (3) & Charge is $+q\left(1-\dfrac1k\right)$\\\hline
(S) Surface D & (4) & Charge is $-q\left(1-\dfrac1k\right)$\\\hline
 & (5) & None\\\hline
\end{tabular}}
\begin{qoptions}
\opt{A}{P $\to$ 4; Q $\to$ 2; R $\to$ 3; S $\to$ 1}
\opt{B}{P $\to$ 2; Q $\to$ 4; R $\to$ 3; S $\to$ 1}
\opt{C}{P $\to$ 1; Q $\to$ 3; R $\to$ 4; S $\to$ 2}
\opt{D}{P $\to$ 2; Q $\to$ 3; R $\to$ 4; S $\to$ 1}
\end{qoptions}
\qn{11} Match the following.
\par\nobreak
{\small\setlength{\tabcolsep}{3pt}\renewcommand{\arraystretch}{1.3}
\begin{tabular}{|L{0.8cm}|L{2.6cm}|L{0.8cm}|L{2.4cm}|}\hline
\multicolumn{2}{|c|}{\textbf{List-I}} & \multicolumn{2}{c|}{\textbf{List-II}}\\\hline
(P) & Spontaneous radioactive decay of an Uranium nucleus initially at rest as given by reaction $^{238}_{92}$U $\to$ $^{234}_{90}$Th $+\,^{4}_{2}$He$+\ldots$ & (1) & Number of proton is increased\\\hline
(Q) & Fusion reaction of two hydrogen nuclei as given by reaction $^{1}_{1}$H $+\,^{1}_{1}$H $\to$ $^{2}_{1}$H$+\ldots$ & (2) & Momentum is conserved\\\hline
(R) & Fission of U$^{235}$ nucleus initiated by a thermal neutron as given by reaction $^{1}_{0}$n $+\,^{235}_{92}$U $\to$ $^{144}_{56}$Ba $+\,^{89}_{36}$Kr $+\,3\,^{1}_{0}$n$+\ldots$ & (3) & Mass is converted to energy or vice versa\\\hline
(S) & $\beta^-$ decay (negative beta decay) & (4) & Charge is conserved\\\hline
 & & (5) & Binding energy is increased\\\hline
\end{tabular}}
\begin{qoptions}
\opt{A}{P $\to$ 2,3,4,5; Q $\to$ 2,3,4,5; R $\to$ 2,3,4; S $\to$ 1,2,3,4,5}
\opt{B}{P $\to$ 2,3,4; Q $\to$ 3,4,5; R $\to$ 2,4,5; S $\to$ 1,2,3,5}
\opt{C}{P $\to$ 2,3; Q $\to$ 2,4,5; R $\to$ 3,5; S $\to$ 1,2,3,4,5}
\opt{D}{P $\to$ 2,3,4,5; Q $\to$ 2,3,4,5; R $\to$ 2,3,4,5; S $\to$ 1,2,3,4,5}
\end{qoptions}
\sect{SECTION-II}
\qn{1} The student assembled an ``almost infinite'' electrical circuit consisting of a very large number of resistors (the diagram of this circuit is shown in the figure). If you move along this circuit from left to right, then after each branching, the resistance of the resistor located above the branching point is equal to the resistance of the resistor located to the left of the branching point, and the resistance of the resistor located below the branching point is twice as large as the resistance of the resistor located above the branching point. The source voltage in this circuit is $U=270$ V. Find the current flowing through the cross-section of the wire at point T if $R=20$ Ohms.
\qfig{p1_phy_s4_q01_1}
\qn{2} A solid sphere of 1 m diameter is kept in an evacuated, enclosure as shown in the diagram. The view factor from the enclosure to the sphere is $\alpha$. The value of $\left(\dfrac{108\alpha}{\pi}\right)$ is
\qfig{p1_phy_s4_q02_1}
\qn{3} Three kg of ice at $-3^\circ$C is exposed to the atmosphere which is at $27^\circ$C. The ice melts and comes into thermal equilibrium with the atmosphere. $c_P$ of ice is 2.0 kJ/kg $^\circ$C and the latent heat of fusion of ice is 333.3 kJ/kg. The change in entropy of ice (in kJ/k) is :- (Rounded off to closest integer) (Take : $\ell$n $(1+x)=x$ for $x\ll1$)
\qn{4} The aphelion distance for Saturn is $1510\times10^6$ km; its perihelion distance is $1350\times10^6$ km. By Kepler's First Law, the Sun is at one focus of this ellipse. The distance of Sun from the other focus is $x\times10^6$ km. The value of $x$ is.
\qn{5} Force acting on a uniform rod having length $3\ell$ \& area of cross section $A$ and Young's modulus $Y$ are shown in figure. find elongation in the rod (take : $\dfrac{F\ell}{YA}=\dfrac43$)
\qfig{p1_phy_s4_q05_1}
\qn{6} A turbine operates under steady flow conditions, receiving steam at the following state : Pressure 1.2 MPa, enthalpy 3000 kJ/kg, velocity 30 m/s and elevation 3 m. The steam leaves the turbine at the following state: Pressure 20 kPa, enthalpy 2500 kJ/kg, velocity 100 m/s, and elevation 0 m. Heat is lost to the surroundings at the rate of 0.3 kJ/s. If the rate of steam flow through the turbine is 0.4 kg/s, what is the power output of the turbine in kW? (Enthalpy is equal to total content of heat of the system equivalent to the system's internal energy plus the product of volume and pressure.) (Round off to nearest integer) ($g=10$ m/s$^2$)
\partsec{2}{CHEMISTRY}
\sect{SECTION-I (i)}
\qn{1} In an analytical determination of arsenic, a solution containing arsenious acid, H$_3$AsO$_3$, KI and a small amount of starch is electrolysed. The electrolysis produces free I$_2$ from I$^-$ ions and the I$_2$ immediately oxidizes the arsenious acid to hydrogen arsenate ion, HAsO$_4^{2-}$.
\[\mathrm{I_2+H_3AsO_3+H_2O\to 2I^-+HAsO_4^{2-}+4H^+}\]
When the oxidation of arsenic is complete, the free iodine combines with the starch to give a deep blue colour. If during a particular run, it takes 96.5 s for a current of 1.68 mA to give an end point (indicated by the blue colour), how many grams of arsenic are present in the solution. [Atomic mass : As $=75$]
\begin{qoptions}
\opt{A}{$6.3\times10^{-4}$ g}
\opt{B}{$1.26\times10^{-4}$ g}
\opt{C}{$3.15\times10^{-5}$ g}
\opt{D}{$6.3\times10^{-5}$ g}
\end{qoptions}
\qn{2} Choose the correct option\\
(1) Ce(IV) can oxidise water and is a good analytical agent\\
(2) Ce(IV) and Yb(II) are diamagnetic in nature and are very strong oxidising agents\\
(3) Mischmetall contains roughly 95\% lanthanoid metal and roughly 5\% magnesium metal\\
(4) Europium has largest atomic radii among the lanthanoids
\begin{qoptions}
\opt{A}{1, 2 are correct statements}
\opt{B}{2, 3 are incorrect statements}
\opt{C}{3, 4 are correct statements}
\opt{D}{1, 4 are incorrect statements}
\end{qoptions}
\qn{3} Incorrect reaction sequence is :
\begin{qoptions}
\opt{A}{\qfigi{p1_chem_s1_q03_A}}
\opt{B}{\qfigi{p1_chem_s1_q03_B}}
\opt{C}{\qfigi{p1_chem_s1_q03_C}}
\opt{D}{CH$_3$--CH=CH--COOH $\xrightarrow[\text{Ni}]{H_2}$ (P) $\xrightarrow[\text{CaO}/\Delta]{\text{NaOH}}$ CH$_3$CH$_3$}
\end{qoptions}
\qn{4} In Nylon-6 which bond among following is absent :
\begin{qoptions}
\opt{A}{Carbon-carbon covalent}
\opt{B}{Amide}
\opt{C}{Ester}
\opt{D}{Carbon hydrogen}
\end{qoptions}
\sect{SECTION-I (ii)}
\qn{5} Standard electrode potential data are useful for understanding the suitability of an oxidant in a redox titration. Some half cell reactions and their standard potentials are given below :\\
MnO$_4^-$(aq.) + 8H$^+$(aq.) + 5e$^-$ $\to$ Mn$^{2+}$(aq.) + 4H$_2$O($\ell$) \hfill $E^0=1.51$ V\\
Cr$_2$O$_7^{2-}$(aq.) + 14H$^+$(aq.) + 6e$^-$ $\to$ 2Cr$^{3+}$(aq.) + 7H$_2$O($\ell$) \hfill $E^0=1.38$ V\\
Fe$^{3+}$(aq.) + e$^-$ $\to$ Fe$^{2+}$(aq.) \hfill $E^0=0.77$ V\\
Cl$_2$(g) + 2e$^-$ $\to$ 2Cl$^-$(aq.) \hfill $E^0=1.40$ V\\
Identify the correct statement(s) regarding the quantitative estimation of aqueous Fe(NO$_3$)$_2$ :
\begin{qoptions}
\opt{A}{MnO$_4^-$ can be used in aqueous HCl.}
\opt{B}{Cr$_2$O$_7^{2-}$ can be used in aqueous HCl.}
\opt{C}{MnO$_4^-$ can be used in aqueous H$_2$SO$_4$.}
\opt{D}{Cr$_2$O$_7^{2-}$ can be used in aqueous H$_2$SO$_4$.}
\end{qoptions}
\qn{6} Gold sol can be prepared by :
\begin{qoptions}
\opt{A}{Peptization}
\opt{B}{By hydrolysis}
\opt{C}{Bredic's arc method}
\opt{D}{Reduction}
\end{qoptions}
\qn{7}
\qfig{p1_chem_s2_q07_1}
\begin{qoptions}
\opt{A}{Ore used contains two metals and is magnetic in nature}
\opt{B}{The same gas (X) is evolved when KClO$_3$(s) is heated}
\opt{C}{Coloured solution (Y) on treatment with hydrogen peroxide in acidic medium, produces a blue colour compound}
\opt{D}{Solution obtained by reaction of W$_{(g)}$ + NaOH$_{(aq.)}$ gives green coloration with acetaldehyde in acidic solution}
\end{qoptions}
\sect{SECTION-I (iii)}
\qn{8} For the reaction\\
2A(g) $\to$ B(g), starting with 4 moles of A in a 5 litre container at 400 K.\\
For all different possible orders of reaction the initial half-life is 40 seconds. Now match following :
{\small\setlength{\tabcolsep}{3pt}\renewcommand{\arraystretch}{1.4}
\begin{center}\begin{tabular}{|L{0.8cm}|L{3.2cm}|L{0.8cm}|L{1.7cm}|}\hline
\multicolumn{2}{|c|}{\textbf{List-I}} & \multicolumn{2}{c|}{\textbf{List-II}}\\\hline
(P) & If reaction follows first order kinetics then initial rate of reaction & (1) & $3.46\times10^{-4}$ M/s\\\hline
(Q) & If reaction follows zero order kinetics then rate of reaction at time : t = 40 sec. & (2) & $6.93\times10^{-3}$ M/s\\\hline
(R) & If reaction follows first order kinetics then rate of reaction at 95\% completion. & (3) & $5\times10^{-3}$ M/s\\\hline
(S) & If reaction follows second order kinetics then rate of reaction at 99\% completion & (4) & $10^{-6}$ M/s\\\hline
 & & (5) & $6.25\times10^{-5}$ M/s\\\hline
\end{tabular}\end{center}}
\begin{qoptions}
\opt{A}{P $\to$ 2; Q $\to$ 3; R $\to$ 1; S $\to$ 4}
\opt{B}{P $\to$ 1; Q $\to$ 2; R $\to$ 3; S $\to$ 4}
\opt{C}{P $\to$ 5; Q $\to$ 1; R $\to$ 2; S $\to$ 4}
\opt{D}{P $\to$ 2; Q $\to$ 1; R $\to$ 4; S $\to$ 3}
\end{qoptions}
\qn{9} Choose the correct option
{\small\setlength{\tabcolsep}{3pt}\renewcommand{\arraystretch}{1.4}
\begin{center}\begin{tabular}{|L{0.8cm}|L{3.1cm}|L{0.8cm}|L{1.9cm}|}\hline
\multicolumn{2}{|c|}{\begin{tabular}[c]{@{}c@{}}\textbf{List-I}\\\textbf{(Reactions)}\end{tabular}} & \multicolumn{2}{c|}{\begin{tabular}[c]{@{}c@{}}\textbf{List-II}\\\textbf{(Observations)}\end{tabular}}\\\hline
(P) & Ferric chloride + Potassium sulphocyanide on reaction produces the following observation & (1) & [Fe(SCN)]$^{2+}$\\\hline
(Q) & CoS + aqua regia, after this, solution is treated with potassium nitrite after neutralisation with NH$_4$OH and the solution is acidified with CH$_3$COOH produces the following observations in the process & (2) & NO gas\\\hline
(R) & Lead sulphate on reaction with ammonium acetate forms a soluble complex and this complex on reaction with potassium iodide produces the following observation & (3) & Yellow ppt.\\\hline
(S) & Deep red fumes turns colourless KOH solution to yellow colour solution, On treating this solution with dil. sulphuric acid and H$_2$O$_2$ + pyridine produces the following observation & (4) & Brown ppt.\\\hline
 & & (5) & Blue ppt.\\\hline
\end{tabular}\end{center}}
\begin{qoptions}
\opt{A}{P $\to$ 1,5; Q $\to$ 2,5; R $\to$ 3; S $\to$ 1}
\opt{B}{P $\to$ 2,5; Q $\to$ 3,4; R $\to$ 5; S $\to$ 3}
\opt{C}{P $\to$ 1; Q $\to$ 2,3; R $\to$ 3; S $\to$ 5}
\opt{D}{P $\to$ 1; Q $\to$ 3,5; R $\to$ 4; S $\to$ 5}
\end{qoptions}
\fwstart
\qn{10} \par\nobreak
\begin{center}\begin{tabular}{|L{0.8cm}|L{4cm}|L{0.8cm}|L{4cm}|}\hline
\multicolumn{2}{|c|}{\begin{tabular}[c]{@{}c@{}}\textbf{List-I}\\\textbf{(Compound)}\end{tabular}} & \multicolumn{2}{c|}{\begin{tabular}[c]{@{}c@{}}\textbf{List-II}\\\textbf{(Statement)}\end{tabular}}\\\hline
(P) & \qfigi{p1_chem_s3_q09_1} & (1) & Is more acidic than carbonic acid\\\hline
(Q) & \qfigi{p1_chem_s3_q09_2} & (2) & Will release CO$_2$ gas with aq. NaHCO$_3$\\\hline
(R) & \qfigi{p1_chem_s3_q09_3} & (3) & Is soluble in aq. NaOH\\\hline
(S) & \qfigi{p1_chem_s3_q09_4} & (4) & Is more acidic than picric acid\\\hline
 & & (5) & Will release H$_2$ gas with Na metal\\\hline
\end{tabular}\end{center}
The correct code is :
\begin{qoptions}
\opt{A}{P $\to$ 2,3; Q $\to$ 4,5; R $\to$ 1,4; S $\to$ 2,3,5}
\opt{B}{P $\to$ 3,4; Q $\to$ 2,3,4,5; R $\to$ 1,2; S $\to$ 2,3}
\opt{C}{P $\to$ 1,2,3,5; Q $\to$ 3,5; R $\to$ 5; S $\to$ 1,2,3,4,5}
\opt{D}{P $\to$ 3,4,5; Q $\to$ 1,2,5; R $\to$ 2,3,4; S $\to$ 1,3,5}
\end{qoptions}
\qn{11} \par\nobreak
\begin{center}\begin{tabular}{|L{0.8cm}|L{5cm}|L{0.8cm}|L{4cm}|}\hline
\multicolumn{2}{|c|}{\begin{tabular}[c]{@{}c@{}}\textbf{List-I}\\\textbf{(For the given reaction)}\end{tabular}} & \multicolumn{2}{c|}{\begin{tabular}[c]{@{}c@{}}\textbf{List-II}\\\textbf{(For final product or product mixture)}\end{tabular}}\\\hline
(P) & \qfigi{p1_chem_s3_q10_1} & (1) & Optically active product mixture\\\hline
(Q) & \qfigi{p1_chem_s3_q10_2} & (2) & Optically inactive product / product mixture\\\hline
(R) & \qfigi{p1_chem_s3_q10_3} & (3) & Product can exhibit optical isomerism\\\hline
(S) & \qfigi{p1_chem_s3_q10_4} & (4) & Carbocation is involved in the reaction sequence\\\hline
 & & (5) & Racemic mixture is obtained\\\hline
\end{tabular}\end{center}
\begin{qoptions}
\opt{A}{P $\to$ 2,4; Q $\to$ 1,3; R $\to$ 2,3,4; S $\to$ 2,3,4,5}
\opt{B}{P $\to$ 1,2,3; Q $\to$ 1,3,4; R $\to$ 1,2; S $\to$ 4,5}
\opt{C}{P $\to$ 4,5; Q $\to$ 1,2,3; R $\to$ 3,4; S $\to$ 2,4}
\opt{D}{P $\to$ 1,2,3; Q $\to$ 1,4; R $\to$ 2,5; S $\to$ 1,2,3}
\end{qoptions}
\fwend
\sect{SECTION-II}
\qn{1} Silver ion forms a complex ion with thiosulphate ion. [Ag(S$_2$O$_3$)$_2$]$^{3-}$, with K$_\mathrm{f}$ = $2.8\times10^{13}$. How many moles of AgBr (K$_\mathrm{sp}$ = $5.4\times10^{-13}$), will dissolve in 100 litre of 0.20 M Na$_2$S$_2$O$_3$ solution. ($\sqrt{15.12}$ = 3.88)\\
(Round off your answers to nearest integer)
\qn{2} You need to convert 1.0 mole of X into product by reaction that follows first order kinetics. At 1.00 p.m. you start the reaction at 27$^\circ$C with one mole of X. At 2.0 p.m. you find that 0.8 mole of X remains. You immediately increase the temperature to 37$^\circ$C. At 3.0 p.m. you discover that 0.4 mole of X are still remaining. You want to be able to finish the reaction by 4.0 p.m. but you cannot stop the reaction until only 0.1 mole of X remains. You decided to increase the temperature again. What is the minimum temperature (in K) required to complete the reaction by 4.0 p.m. ? [Use : ln 2 = 0.7, ln 5 = 1.6, ln 7 = 1.9, R = 8J/K-mole]\\
(Round off your answer to nearest integer)
\qn{3} Number of correct statement(s) :-\\
(1) Zn , Cd , Hg and Mn have typical metallic structure at normal temperature\\
(2) Zn has the lowest enthalpy of atomization among the 3d series metals\\
(3) In general density of 3d series transition metals increases as we move more from left to right in a period\\
(4) Cr$^{+2}$ is reducing and Mn$^{+3}$ is oxidizing in nature\\
(5) Cu$^{+1}$ compounds in aqueous solution undergoes disproportionation reaction\\
(6) Ferrates readily decomposed into ferric oxide and O$_2$(g)\\
(7) In general tendency to form divalent cations across the 3d series decreases because the sum of the IE$_1$ + IE$_2$ increases as we move from left to right in a period\\
(8) Interstitial compounds are chemically inert and they retain their metallic conductivity
\qn{4} `X' are the number of complexes which are stable than [Co(CN)$_6$]$^{3-}$ among the following\\
{[}CoCl$_6$]$^{4-}$ , [Co(NH$_3$)$_6$]$^{3+}$ , [Co(H$_2$O)$_6$]$^{3+}$ , [Co(ox)$_3$]$^{3-}$ , [Co(en)$_3$]$^{3+}$\\
`Y' are number of linkage isomers of [Co(SCN)$_4$(NO$_2$)$_2$]$^{3-}$\\
`Z' are number of compounds which show geometrical isomerism among the following compounds\\
{[}Pt(en)$_2$]$^{2+}$ , [Pt(en)Cl$_2$]$^{2+}$ , [Pt(gly)Cl$_2$] , [Pt(bn)$_2$]$^{2+}$\\
Find $\dfrac{\mathrm{X}+\mathrm{Y}}{\mathrm{Z}}$.
\qn{5} (i) Boric acid on heating yields a compound Y , Y on heating yields a compound Z, the oxidation state of Z is `A'.\\
(ii) Metal Borohydrides are prepared by the reaction of metal hydrides with compound `C' in diethyl ether, number of hybrid orbitals involves in one molecule of product formed is `B'.\\
Find A + B
\qn{6} Number of reactions in which aldehyde is obtained by oxidation of organic reactant ?
\qfig{p1_chem_s4_q06_1}
\partsec{3}{MATHEMATICS}
\sect{SECTION-I (i)}
\qn{1} The graph of $2x^2 + xy + 3y^2 - 11x - 20y + 40 = 0$ is an ellipse in the first quadrant of the xy-plane. Let a and b be the maximum and minimum values of $\dfrac{y}{x}$ over all points (x, y) on the ellipse. What is the value of a + b ?
\begin{qoptions}
\opt{A}{3}
\opt{B}{$\sqrt{10}$}
\opt{C}{$\dfrac{7}{2}$}
\opt{D}{$\dfrac{9}{2}$}
\end{qoptions}
\qn{2} How many 15-letter arrangements of 5 A's, 5 B's and 5 C's have no A's in the first 5 letters, no B's in the next 5 letters and no C's in the last 5 letters ?
\begin{qoptions}
\opt{A}{$\displaystyle\sum_{k=0}^{5}\binom{5}{k}^{3}$}
\opt{B}{$3^5\cdot 2^5$}
\opt{C}{$2^{15}$}
\opt{D}{$\dfrac{15!}{(5!)^3}$}
\end{qoptions}
\qn{3} If $f(x) = 3\left\{\dfrac{[x]}{3}\right\}$, $g(x) = \left[\dfrac{1}{1+x^2}\right]$ and $\phi(x) = g(f(x))\{x\}^2 + g(f(x)-1) + g(f(x)-2)(1-\{x\})$ where [x] and \{x\} denote integral and fractional parts of x respectively, then $\displaystyle\int_0^{2024}\phi(x)\,dx =$
\begin{qoptions}
\opt{A}{1235}
\opt{B}{1233}
\opt{C}{$1227\dfrac{1}{3}$}
\opt{D}{$1228\dfrac{1}{3}$}
\end{qoptions}
\qn{4} The image of the point (3, 5) in the line $x - y + 1 = 0$, lies on :
\begin{qoptions}
\opt{A}{$(x-2)^2 + (y-2)^2 = 12$}
\opt{B}{$(x-2)^2 + (y-4)^2 = 4$}
\opt{C}{$(x-4)^2 + (y+2)^2 = 16$}
\opt{D}{$(x-4)^2 + (y-4)^2 = 8$}
\end{qoptions}
\sect{SECTION-I (ii)}
\qn{5} Which of the following options are correct ?
\begin{qoptions}
\opt{A}{If $a \ge b > 1$, then the largest possible value of $\log_a\left(\dfrac{a}{b}\right) + \log_a\left(\dfrac{b}{a}\right)$ is 0}
\opt{B}{Number of integers `b' for which $\log_b 2024$ is a positive integer is 0.}
\opt{C}{There is a unique positive integer `n' such that $\log_2(\log_{16}n) = \log_4(\log_4 n)$. The sum of the digits of `n' is 13.}
\opt{D}{The set of real numbers `x' for which $\log_{2024}\log_{2023}\log_{2022}\log_{2021}\log_{2020}x$ is defined is $\{x|x>c\}$. Then `c' is $2020^{2021^{2022}}$}
\end{qoptions}
\qn{6} Let f(x) and g(x) are differentiable functions such that $f(x) + \displaystyle\int_0^x g(t)\,dt = \sin x(\cos x - \sin x)$, and $(f'(x))^2 + (g(x))^2 = 1$ then f(x) and g(x) respectively, are
\begin{qoptions}
\opt{A}{$\cos^2x$, $\cos 2x$}
\opt{B}{$\dfrac{\cos 2x}{2}$, $\cos 2x$}
\opt{C}{$\dfrac{1}{2}\sin 2x$, $-\sin 2x$}
\opt{D}{$-\sin^2x$, $\cos 2x$}
\end{qoptions}
\qn{7} If f(x) be such that $f(x) = \max(|3-x|,\ 3-x^3)$ then :
\begin{qoptions}
\opt{A}{f(x) is continuous $\forall$ x $\in$ R}
\opt{B}{f(x) is derivable at three points only}
\opt{C}{f(x) is non-derivable at three points only}
\opt{D}{f(x) is non-derivable at four points only}
\end{qoptions}
\sect{SECTION-I (iii)}
\qn{8} Match the following from List-I to List-II.
{\small\setlength{\tabcolsep}{3pt}\renewcommand{\arraystretch}{1.5}
\begin{center}\begin{tabular}{|L{0.8cm}|L{3.2cm}|L{0.8cm}|L{1.0cm}|}\hline
\multicolumn{2}{|c|}{\textbf{List-I}} & \multicolumn{2}{c|}{\textbf{List-II}}\\\hline
(P) & The equation of tangent of the ellipse $\dfrac{x^2}{25}+\dfrac{y^2}{16}=1$ which cuts off numerically equal intercepts on axes is $x-y=a$ where $|a|$ equal to & (1) & $\sqrt2$\\\hline
(Q) & The normal $y = mx - 2am - am^3$ to the parabola $y^2=4ax$ subtends a right angle at the vertex if $|m|$ equal to & (2) & $\sqrt3$\\\hline
(R) & The equation of the common tangent to parabola $y^2=4ax$ and $x^2=4y$ is $x+y+\dfrac{n}{\sqrt3}=0$, then n is equal to & (3) & $\sqrt8$\\\hline
(S) & If the equation of common tangent to parabola $y^2=8x$ and the hyperbola $3x^2-y^2=3$ is $4x-2y+\dfrac{\lambda}{\sqrt2}=0$, then $|\lambda|$ is equal to & (4) & $\sqrt{41}$\\\hline
 & & (5) & 2\\\hline
\end{tabular}\end{center}}
\begin{qoptions}
\opt{A}{P $\to$ 1; Q $\to$ 2; R $\to$ 3; S $\to$ 4}
\opt{B}{P $\to$ 2; Q $\to$ 4; R $\to$ 1; S $\to$ 3}
\opt{C}{P $\to$ 2; Q $\to$ 1; R $\to$ 4; S $\to$ 3}
\opt{D}{P $\to$ 4; Q $\to$ 1; R $\to$ 2; S $\to$ 3}
\end{qoptions}
\fwstart
\qn{9} Match the following.\par\nobreak
\begin{center}\begin{tabular}{|L{0.8cm}|L{10cm}|L{0.8cm}|L{1.2cm}|}\hline
\multicolumn{2}{|c|}{\textbf{List-I}} & \multicolumn{2}{c|}{\textbf{List-II}}\\\hline
(P) & $(2x-10y^3)\dfrac{dy}{dx}+y=0$, $y(1)=1$ then the value of $[x]$ when $y=2$ is equal to ([.] GIF) & (1) & 3\\\hline
(Q) & $(x^2y^2-1)dy+2xy^3dx=0$, $y(1)=1$ then at $y=2$ the value $4x^2$ equals & (2) & 7\\\hline
(R) & Area of the region formed by the points $(x, y)$ satisfying $[x]^2=[y]^2$ for $0\le x\le 4$ is greater than ([.] GIF) & (3) & 15\\\hline
(S) & Let the function $f(x)$ which is continuous every where, and satisfies $f(x)+\dfrac{1}{f(x)}=e^x+e^{-x}$ then the area bounded by $y=f(x)$, x-axis, $x=-1$ and $x=1$ can have integral part & (4) & 2\\\hline
 & & (5) & 1\\\hline
\end{tabular}\end{center}
\begin{qoptions}
\opt{A}{P $\to$ 3; Q $\to$ 1; R $\to$ 1,4,5; S $\to$ 4}
\opt{B}{P $\to$ 2; Q $\to$ 3; R $\to$ 2,3,4; S $\to$ 5}
\opt{C}{P $\to$ 5; Q $\to$ 1,2,3; R $\to$ 4; S $\to$ 2}
\opt{D}{P $\to$ 3,4,5; Q $\to$ 2; R $\to$ 1; S $\to$ 3}
\end{qoptions}
\qn{10} Match List-I with List-II and select the correct answer using the code given below the list.\par\nobreak
\begin{center}\begin{tabular}{|L{0.8cm}|L{10cm}|L{0.8cm}|L{1.2cm}|}\hline
\multicolumn{2}{|c|}{\textbf{List-I}} & \multicolumn{2}{c|}{\textbf{List-II}}\\\hline
(P) & Number of real roots of $2^{-x}+2^{-x^3}-x=0$ is & (1) & 2\\\hline
(Q) & Let $\theta_1+\theta_2=2\theta$ (where $\theta$ is acute angle) then $\dfrac{2(\sin\theta_1-\sin\theta_2)}{(\cos\theta_2-\cos\theta_1)}\tan\theta$ is (wherever defined) & (2) & 3\\\hline
(R) & $6\left(\log_{20}\left(\sqrt{6-2\sqrt5}+\sqrt{6+2\sqrt5}\right)\right)$ is & (3) & 4\\\hline
(S) & Number of integers satisfying the inequality $\log_{\frac1e}|x-4|>-1$ is & (4) & 1\\\hline
 & & (5) & 6\\\hline
\end{tabular}\end{center}
The correct option is :
\begin{qoptions}
\opt{A}{P $\to$ 4; Q $\to$ 1; R $\to$ 2; S $\to$ 3}
\opt{B}{P $\to$ 1; Q $\to$ 2; R $\to$ 3; S $\to$ 4}
\opt{C}{P $\to$ 4; Q $\to$ 2; R $\to$ 3; S $\to$ 1}
\opt{D}{P $\to$ 1; Q $\to$ 4; R $\to$ 3; S $\to$ 2}
\end{qoptions}
\qn{11} Match the following\par\nobreak
\begin{center}\begin{tabular}{|L{0.8cm}|L{10cm}|L{0.8cm}|L{1.2cm}|}\hline
\multicolumn{2}{|c|}{\textbf{List-I}} & \multicolumn{2}{c|}{\textbf{List-II}}\\\hline
(P) & The value of $\displaystyle\int_0^1\sqrt[5]{2x^3-3x^2+3x-1}\,dx$ is less than & (1) & 0\\\hline
(Q) & Let $I_1=\displaystyle\int_0^{\pi/2}\cos\left(\dfrac{2\pi}{3}\sin^2x\right)dx$, $I_2=\displaystyle\int_0^{\pi/2}\cos\left(\dfrac{\pi}{3}\sin x\right)dx$ then the value of $\dfrac{I_2}{I_1}$ is greater than & (2) & 1\\\hline
(R) & $\displaystyle\int x(x+1)\left(\dfrac{1}{x^8}+\dfrac{2}{x^9}\right)^{1/10}dx=\phi(x)$, $\phi(0)=0$, $\phi(1)=\dfrac{a}{11}(3)^{\frac{11}{b}}$ then $\dfrac{b}{a}$ is less than & (3) & 2\\\hline
(S) & Let $f(x)=\displaystyle\int x^2(2x+6\tan x-2x\tan^2x)\cos^2x\,dx$ and $y=f(x)$ passes through $(\pi, 0)$ then the number of solutions of $f(x)=x^3$ is less than & (4) & 3\\\hline
 & & (5) & 4\\\hline
\end{tabular}\end{center}
\begin{qoptions}
\opt{A}{P $\to$ 1; Q $\to$ 2,3; R $\to$ 3,4,5; S $\to$ 1,2}
\opt{B}{P $\to$ 2,3,4,5; Q $\to$ 1,2; R $\to$ 4,5; S $\to$ 5}
\opt{C}{P $\to$ 3; Q $\to$ 2,4,5; R $\to$ 3,4; S $\to$ 1,2}
\opt{D}{P $\to$ 1,2; Q $\to$ 1,3,5; R $\to$ 4; S $\to$ 1,5}
\end{qoptions}
\fwend
\sect{SECTION-II}
\qn{1} Two distinct, real, infinite geometric series each have a sum of 1 and have the same second term. The third term of one of the series is 1/8, and the second term of both series can be written in the form $\dfrac{\sqrt{m}-n}{p}$ where m, n and p are positive integers and m is not divisible by the square of any prime. Find $100m + 10n + p$.
\qn{2} Let P(x) be a quadratic polynomial with real coefficients satisfying $x^2 - 2x + 2 \le P(x) \le 2x^2 - 4x + 3$ for all real numbers x, and suppose P(11) = 181. Find P(16).
\qn{3} Find
{\small\[L=\lim_{x\to\infty}\left(\sqrt{x}\sin\left(\sin\left(\sin\left(\lim_{n\to\infty}\sum_{k=1}^{n}\frac{1}{\sqrt{xn^2+k}}\right)\right)\right)\right)\]}
\qn{4} Let $A_n = [a_{ij}]$ be a square matrix of order 3, where $a_{ij} = \dfrac{2i+j}{3^{2n}}$ for all i and j, $1\le i, j\le 3$. If $\ell = \displaystyle\lim_{n\to\infty}\mathrm{Tr.}(3A_1 + 3^2A_2 + 3^3A_3 + \ldots.. 3^nA_n)$, then `$\ell$' is equal to\\
(where Tr.A means trace of A)
\qn{5} If the ratio of the 4$^\mathrm{th}$ term to the 2$^\mathrm{nd}$ term in the expansion of $\left(\dfrac{1}{x}+x\tan x\right)^5$ is $\dfrac{2}{27}\pi^4$ then the value of [x] \underline{\hspace{1.5cm}}\\
([.] denotes G.l.F.) where $x>0$
\qn{6} Let $f(x) = x^\alpha + 3x^\beta$ with $\alpha>1$, $\beta>1$. If the value of a for which the area of the region bounded by $y = f(x)$ and straight lines $x = 0$, $x = 1$ and $y = f(a)$ is minimum is $\lambda$, then the value of $70\lambda$ is
\papersec{2}
\partsec{1}{PHYSICS}
\sect{SECTION-I (i)}
\qn{1} A uniform ring of mass m, radius a, resistance R, modulus of elasticity Y, area of cross section A is kept on a smooth surface. When it is somehow expanded \& released, it undergoes radial oscillations with angular frequency $\omega_0$. If the same ring is made to oscillate radially in presence of vertical magnetic field B, it undergoes damped oscillation with angular frequency $\omega$.
\qfig{p2_phy_s1_q01_1}
\begin{qoptions}
\opt{A}{$\omega_0=\sqrt{\dfrac{4\pi YA}{am}}$}
\opt{B}{$\omega=\sqrt{\dfrac{4\pi YA}{ma}-\dfrac{2\pi^4a^4B^4}{m^2R^2}}$}
\opt{C}{$\omega_0=\sqrt{\dfrac{2\pi YA}{am}}$}
\opt{D}{$\omega=\sqrt{\dfrac{2\pi YA}{ma}-\dfrac{2\pi^4a^4B^4}{m^2R^2}}$}
\end{qoptions}
\qn{2} The pitch of screw gauge is 1 mm and its circular scale is divided into 100 divisions. When nothing is put between the studs the circular scale lies between 2$^\mathrm{nd}$ and 3$^\mathrm{rd}$ division of main scale, but when circular scale is rotated by 270$^\circ$ the 3$^\mathrm{rd}$ division of main scale is seen and the 3$^\mathrm{rd}$ division of main scale coincide with the 12$^\mathrm{th}$ division of circular scale. When a glass plate is placed between the studs, the circular scale lies between 28$^\mathrm{th}$ to 29$^\mathrm{th}$ division of main scale and circular scale reads 20 divisions. Find the thickness of the glass plate in millimeter. (Assume value of one main scale division is same as pitch)
\begin{qoptions}
\opt{A}{25.83 mm}
\opt{B}{26.83 mm}
\opt{C}{25.73 mm}
\opt{D}{25.80 mm}
\end{qoptions}
\qn{3} A solid cone is kept in three immiscible liquids of density $\rho$, $2\rho$ and $3\rho$ inside a container. The radius of cone is `R' and it's height is 4R. Container is accelerated with a constant acceleration g/2 in vertically upward direction as shown. If cone is in equilibrium relative to the frame of container then density of the cone is
\qfig{p2_phy_s1_q03_1}
\begin{qoptions}
\opt{A}{$\dfrac{73}{24}\rho$}
\opt{B}{$\dfrac{146}{25}\rho$}
\opt{C}{$\dfrac{39}{16}\rho$}
\opt{D}{$\dfrac{146}{37}\rho$}
\end{qoptions}
\qn{4} Consider four point masses each of mass m are fixed at the vertex of a Rhombus of side $\ell$. The escape speed from the centre of the Rhombus is (G is universal gravitational constant)
\qfig{p2_phy_s1_q04_1}
\begin{qoptions}
\opt{A}{$\sqrt{\dfrac{8Gm}{\ell}}$}
\opt{B}{$\sqrt{4\sqrt2\,\dfrac{Gm}{\ell}}$}
\opt{C}{$\sqrt{\dfrac{8Gm}{\ell}\left(1+\dfrac{1}{\sqrt3}\right)}$}
\opt{D}{$\sqrt{\dfrac{8Gm}{\ell}\left(1+\dfrac{1}{\sqrt5}\right)}$}
\end{qoptions}
\sect{SECTION-I (ii)}
\qn{5} A uniform rod of length 4 m and mass $2\sqrt2$ kg revolves with constant angular velocity $\omega$ about a vertical axis through a smooth joint A at one extremity of the rod so that it describes a cone of semi vertical angle $53^\circ$ as shown in the figure. Choose the correct statement(s). (g = 10 m/s$^2$)
\qfig{p2_phy_s2_q05_1}
\begin{qoptions}
\opt{A}{The value of angular velocity w is 2.5 rad/s}
\opt{B}{The angle between hinge reaction at point A and vertical axis AC is $45^\circ$}
\opt{C}{The magnitude of the hinge reaction at point A is 40 N}
\opt{D}{The angular momentum of the rod is constant}
\end{qoptions}
\qn{6} An imaginary particle has a charge equal to the charge of electron but mass of imaginary particle is 50 times the mass of electron. Now, consider that imaginary particle moves in circular orbit around a nucleus. The charge of nucleus is +4e, and nucleus is assumed to be at rest. Assume mass of nucleus is very large as compared to the mass of imaginary particle moving around it. Choose the correct option(s).
\begin{qoptions}
\opt{A}{The radius of n$^\mathrm{th}$ orbit is $\dfrac{n^2h^2\varepsilon_0}{200\pi me^2}$, where m is mass of electron, h is Planck constant, e is charge of electron.}
\opt{B}{If all atoms are at 16$^\mathrm{th}$ excited state then during de-excitation till 10$^\mathrm{th}$ excited state the total maximum number of possible wavelengths emitted is 21.}
\opt{C}{The wavelength of photon emitted when transition takes place from third excited state to ground state is 1.21 A.}
\opt{D}{Since mass of imaginary particle is different from the mass of electron so the velocity of imaginary particle and the electron in n$^\mathrm{th}$ orbit will be different.}
\end{qoptions}
\qn{7} A narrow slit S transmitting light of wavelength $\lambda$ is placed a distance d above a large plane mirror as shown in the figure. The light coming directly form the slit and that coming after reflection from the plane mirror interfere on a screen $\Sigma$ placed at a distance D from the slit. Assume that the mirror is 100\% reflecting.\\
(Given D = 1m, d = 0.5 mm and $\lambda$ = 500 nm)
\qfig{p2_phy_s2_q07_1}
\begin{qoptions}
\opt{A}{Distance from O, where closest interference maximum is formed on the screen is $\dfrac14$ mm.}
\opt{B}{Distance from O, where closest interference minimum is formed on the screen is $\dfrac14$ mm.}
\opt{C}{Fringe width of the interference pattern formed on the screen is $\dfrac12$ mm.}
\opt{D}{Closest point to O on the screen where waves interfere with a phase difference of $\pi$ is $\dfrac18$ mm.}
\end{qoptions}
\sect{SECTION-II (i)}
\paraq{1}{2}
A hollow cylinder of radius 4R is rotating about fixed horizontal axis passing through point `O' with angular velocity $\omega_0$. A solid cylinder of radius `R' is rolling without slipping with respect to inner surface of hollow cylinder. At the given instant, the line OC has angular velocity of $2\omega_0$. Point A and B are topmost and bottom most point of solid cylinder respectively and C is its centre. Then, at the given instant
\qfig{p2_phy_s3_para1_2_1}
\qn{1} Angular velocity of solid cylinder about its own axis is $\alpha\omega_0$, the value of $\alpha$ is :-
\qn{2} The magnitude of acceleration of point B of solid cylinder with respect to the contact surface of hollow cylinder is $\beta\omega_0^2$R, the value of $\beta$ is :-
\paraq{3}{4}
In the set-up shown `E' is unknown emf with internal resistance `r'. AB is a uniform potentiometer wire, 11m long and resistance $1\Omega\,\mathrm{m}^{-1}$. When AJ = 1m, current through galvanometer is 2.4 A when AJ = 10 m, current through galvanometer is $\dfrac{6}{11}$A, in the opposite direction as in previous case. Galvanometer has negligible resistance.
\qfig{p2_phy_s3_para3_4_1}
\qn{3} The value of E is (in V)
\qn{4} The value of r is (in $\Omega$)
\columnbreak
\sect{SECTION-II (ii)}
\qn{5} A particle of mass m = 9 gm is trapped between two perfectly rigid parallel walls. The particle bounces back and forth between the walls without losing any energy. From a wave point of view, the particle trapped between the walls is like a standing wave in stretched string between the walls. Distance between the two walls is L = 1 m. Calculate the energy difference between third energy state and the ground state (lowest energy state) of the particle. Answer the energy difference in C $\times$ $10^{-65}$ Joule. (Take Planck's constant h = $6\times10^{-34}$ J-sec)
\qn{6} Two short electric dipoles each of dipole moment P are placed at origin with direction along the coordinate axis as shown in the figure. Calculate the work done (in joule) in moving a point charge $q_0$ from coordinate (3, $-$4) to (3, 4) ? (Take : $\dfrac{Pq_0}{4\pi\varepsilon_0}$ = 125 SI unit)
\qfig{p2_phy_s4_q06_1}
\qn{7} In a container of depth `9a' a triangular wedge is kept touching a wall as shown such that its one edge is just visible to an external observer. When a liquid of refractive index $\mu$ is completely filled in the container top of wedge also becomes visible. If the refractive index of liquid is $\sqrt{\dfrac{k}{10}}$, the value of k is :-
\qfig{p2_phy_s4_q07_1}
\qn{8} A string of 0.2 kg/m and length 1.2 m is fixed at both ends are stretched such that tension generated in the string is 320 N. The string vibrates in six segments with amplitude of 0.25 cm. Then maximum transverse velocity of a particle which is at a distance of $\dfrac{1}{30}$m from the fixed end is k$\pi$ cm/sec. Write value of k.
\qn{9} A uniform disc of radius R having charge Q distributed uniformly all over its surface is placed on a smooth horizontal surface. A magnetic field, B = $kxt^2$, where k is a constant, x is the distance (in metre) from the centre of the disc and t is the time (in second), is switched on perpendicular to the plane of the disc. Find the torque (in N-m) acting on the disc after 15 sec. (Take 4kQ = 5 S.I. unit and R = 1m).
\qfig{p2_phy_s4_q09_1}
\qn{10} 24 resistors of 7$\Omega$ are arranged in a square as shown. Equivalent resistance of the network between points A and B is :-
\qfig{p2_phy_s4_q10_1}
\partsec{2}{CHEMISTRY}
\sect{SECTION-I (i)}
\qn{1} \qfig{p2_chem_s1_q01_1}
\begin{qoptions}
\opt{A}{\qfigi{p2_chem_s1_q01_A}}
\opt{B}{\qfigi{p2_chem_s1_q01_B}}
\opt{C}{\qfigi{p2_chem_s1_q01_C}}
\opt{D}{\qfigi{p2_chem_s1_q01_D}}
\end{qoptions}
\qn{2} A plot for $\log\dfrac{x}{m}$ vs $\log P$ shows a straight line with an angle $60^\circ$ and intercept of 0.3010 on y-axis. The amount of gas adsorbed per gram of adsorbent if initial pressure maintained is 1 atm is:
\begin{qoptions}
\opt{A}{1 gm}
\opt{B}{2 gm}
\opt{C}{3 gm}
\opt{D}{4 gm}
\end{qoptions}
\qfullbegin
\qn{3} \qfig{p2_chem_s1_q03_1}
\qfullsplit
Select correct option regarding compound `X' and `Y' :
\begin{qoptions}
\opt{A}{`X' has sp$^3$d hybridisation but `Y' has sp$^3$.}
\opt{B}{Both `X' and `Y' have same total number of p$\pi$ $-$ d$\pi$ bonds.}
\opt{C}{Individual hydrolysed solution of `X' and `Y' converts blue litmus paper to red.}
\opt{D}{Both compound are planar.}
\end{qoptions}
\qfullend
\qn{4} Which of the following statement is \textbf{INCORRECT} ?
\begin{qoptions}
\opt{A}{Cassiterite ore of tin contains the impurities of Woll-framite which are separated by magnetic separator.}
\opt{B}{Tin metal is obtained by the carbon reduction of concentrated ore.}
\opt{C}{In the extraction of lead from galena, the roasting and self-reduction are carried in the same furnace at different temperatures.}
\opt{D}{Primary reducing agent of haematite in blast-furnace is coke in upper part and CO in lower part of furnace.}
\end{qoptions}
\sect{SECTION-I (ii)}
\qn{5} In which of the following reaction product formation takes place by Hofmann rule?
\begin{qoptions}
\opt{A}{\qfigi{p2_chem_s2_q05_A}}
\opt{B}{\qfigi{p2_chem_s2_q05_B}}
\opt{C}{\qfigi{p2_chem_s2_q05_C}}
\opt{D}{\qfigi{p2_chem_s2_q05_D}}
\end{qoptions}
\qn{6} \qfig{p2_chem_s2_q06_1}
Identify the amino acids obtained by hydrolysis of the above compound :
\begin{qoptions}
\opt{A}{Glycine}
\opt{B}{Phenylalonine}
\opt{C}{Alanine}
\opt{D}{Glutamic acid}
\end{qoptions}
\qn{7} A d-block metal ion with d$^\mathrm{n}$ configuration, when it forms tetrahedral complex but its spin magnetic moment remains same either with 4H$_2$O or with 4Cl$^-$ ligand, which of the following statements are INCORRECT ?
\begin{qoptions}
\opt{A}{Element always form a colourless compound}
\opt{B}{Number of electrons in e-set of d-orbitals are always higher than in t2-set of d-orbitals}
\opt{C}{It can have d$^5$ or d$^8$ configuration}
\opt{D}{Colour of the complexes must be same}
\end{qoptions}
\sect{SECTION-II (i)}
\paraq{1}{2}
Following is the titration curve of two weak monoprotic acids HA and HB (5 milli-moles each) in 50 ml aqueous solution titrated against aqueous solution of strong base NaOH (0.1 M) in two different experiments : [Given : log 2 = 0.3]
\qfig{p2_chem_s3_para1_2_1}
\qn{1} What is pK$_\mathrm{a}$ value for acid HB ?
\qn{2} pH at equivalence point if HA is titrated by NaOH in the experiment given above :
\columnbreak
\paraq{3}{4}
\qfig{p2_chem_s3_q02_1}
\qn{3} Total number of orbitals used in bonding in anion of compounds X and Y are :
\qn{4} The pink solution and the colourless solution in the above scheme, combined to give a product in presence of ZnO catalyst. The oxidation number of central atom in the metal containing product is :
\sect{SECTION-II (ii)}
\qn{5} The total number of cyclic, structural isomers for a compound with molecular formula C$_5$H$_{10}$ is \underline{\hspace{1cm}}
\qn{6} \qfig{p2_chem_s4_q06_1}
Double bond equivalent of final product is \underline{\hspace{1cm}}
\qn{7} \qfig{p2_chem_s4_q07_1}
How many products are obtained in the above reaction.
\qn{8} A Hydrogen electrode filled with hydrogen gas at a pressure of 1 atm was dipped in a solution of weak monoprotic acid (HA), 20 mL of a NaOH is added and connected to a reference electrode [Ag/AgCl/Cl$^-$ (0.1 M) electrode]. E.m.f. of the cell was found to be 0.47 V. To the resulting solution, 30 mL of NaOH of the same strength was further added and new e.m.f. was found to be 0.53 V. What will be the pKa value for weak monoprotic acid (HA) ?\\
{[}$E^0_{\mathrm{Cl^-/AgCl/Ag}}$ = 0.23V , $2.303\dfrac{RT}{F}$ = 0.06, log 2 = 0.3{]}\\
(Round off your answer to nearest integer)
\qn{9} At what temperature (in Kelvin) is the following reaction at equilibrium, when all substances are at standard pressure? Assuming that entropies and enthalpies of reaction do not vary with temperature.\\
PCl$_{3(\mathrm{g})}$ + Cl$_{2(\mathrm{g})}$ $\rightleftharpoons$ PCl$_{5(\mathrm{g})}$
\par\smallskip
\begin{center}\small\begin{tabular}{|c|c|c|}\hline
 & $\Delta$H$^0_\mathrm{f}$ (kJ/mole) & S$^\circ$ (J/mole-K)\\\hline
PCl$_{3(\mathrm{g})}$ & $-290$ & 310\\\hline
Cl$_{2(\mathrm{g})}$ & 0 & 220\\\hline
PCl$_{5(\mathrm{g})}$ & $-375$ & 360\\\hline
\end{tabular}\end{center}
\qn{10} One mole of an ideal monoatomic gas undergoes two reversible processes (A $\to$ B and B $\to$ C) as shown in the given figure :
\qfig{p2_chem_s4_q10_1}
A $\to$ B is an adiabatic process. If the total heat absorbed in the entire process (A $\to$ B and B $\to$ C) is (RT$_2$ ln 10), then value of (2 log V$_3$) is \underline{\hspace{1.5cm}}.\\
{[}Use, molar heat capacity of the gas at constant pressure, C$_{\mathrm{p,m}}$ = $\dfrac52$ R{]}
\partsec{3}{MATHEMATICS}
\sect{SECTION-I (i)}
\qn{1} If $x_1$, $x_2$ and $x_3$ are the positive roots of the equation $x^3 - 6x^2 + 3px - 2p = 0$, $p \in R - \{0\}$ then the value $\sin^{-1}\left(\dfrac{1}{x_1}+\dfrac{1}{x_2}\right) + \cos^{-1}\left(\dfrac{1}{x_2}+\dfrac{1}{x_3}\right) - \tan^{-1}\left(\dfrac{1}{x_3}+\dfrac{1}{x_1}\right)$ is equal to
\begin{qoptions}
\opt{A}{$\dfrac{\pi}{4}$}
\opt{B}{$\dfrac{\pi}{2}$}
\opt{C}{$\dfrac{3\pi}{4}$}
\opt{D}{$\pi$}
\end{qoptions}
\qn{2} $\displaystyle\sum_{r=0}^{10}2^r\,{}^{20}C_r\,{}^{20-r}C_{10-r}$ is equal to
\begin{qoptions}
\opt{A}{${}^{20}C_{10}$}
\opt{B}{${}^{20}C_{10}\left(\dfrac{3}{2}\right)^{10}$}
\opt{C}{${}^{20}C_{10}\,3^{10}$}
\opt{D}{${}^{20}C_{10}\,2^{10}$}
\end{qoptions}
\qn{3} The solution of $\cos\left(\dfrac{x}{y}\right)(y\,dx - x\,dy) = xy^3(x\,dy + y\,dx)$, is
\begin{qoptions}
\opt{A}{$\cos\left(\dfrac{x}{y}\right)=\dfrac{(xy)^2}{2}+C$}
\opt{B}{$\sin\left(\dfrac{x}{y}\right)=\dfrac{(xy)^2}{2}+C$}
\opt{C}{$\cos\left(\dfrac{x}{y}\right)=\dfrac{x^2y}{2}+C$}
\opt{D}{$\cos\left(\dfrac{x}{y}\right)=\dfrac{xy^2}{2}+C$}
\end{qoptions}
\qn{4} If $\displaystyle\lim_{x\to1}\sec^{-1}\left(\dfrac{k}{\ell nx}-\dfrac{k}{x-1}\right)=\dfrac{\pi}{3}$, then the value of k is equal to
\begin{qoptions}
\opt{A}{2}
\opt{B}{3}
\opt{C}{4}
\opt{D}{6}
\end{qoptions}
\sect{SECTION-I (ii)}
\qn{5} If the equation $\sec x+\mathrm{cosec}\,x=k$, ($k\in R$) has exactly 3 distinct solutions between 0 and $\pi$, then k can take values
\begin{qoptions}
\opt{A}{$\sqrt2$}
\opt{B}{2}
\opt{C}{$3\sqrt2$}
\opt{D}{4}
\end{qoptions}
\qn{6} If the equation $x^4 - 4x^3 + ax^2 + bx + 1 = 0$ has four positive roots then
\begin{qoptions}
\opt{A}{the eccentricity of $\dfrac{x^2}{a^2}+\dfrac{y^2}{b^2}=1$ is $\dfrac{\sqrt5}{3}$}
\opt{B}{the eccentricity of $\dfrac{x^2}{a^2}-\dfrac{y^2}{b^2}=1$ is $\dfrac{\sqrt{13}}{3}$}
\opt{C}{the latus rectum of $y^2=4ax$ is 12}
\opt{D}{the latus rectum of $x^2=4by$ is 24}
\end{qoptions}
\qn{7} The focus of parabola is (2, 3) and it is touching the coordinate axes. Then
\begin{qoptions}
\opt{A}{equation of axis of parabola is $3x-2y=0$}
\opt{B}{length of latus rectum of parabola is $\dfrac{24}{\sqrt{13}}$}
\opt{C}{equation of tangent at vertex is $3x+2y-6=0$}
\opt{D}{equation of directrix is $3x+2y=0$}
\end{qoptions}
\sect{SECTION-II (i)}
\paraq{1}{2}
Consider $f(x)=x^4+ax^3+3x^2+bx+2$ and $g(x)=x^4+2x^3-bx^2+3x+c$, where a, b, c $\in$ R, and $\displaystyle\lim_{x\to\infty}\left(\sqrt{f(x)}-\sqrt{g(x)}\right)=4$.
\qn{1} If a, b, c are three consecutive terms of HP. then c equals
\qn{2} If $f(x)-g(x)>1$, $\forall$ x $\in$ R then number of positive integral value of c is
\paraq{3}{4}
10 digit numbers are formed by using all the digits 0, 1, 2, ..., 9 such that they are divisible by 11111. Then, answer the following questions.
\qn{3} The digit in tens place, in smallest of such numbers is
\qn{4} The digit in units place, in greatest of such numbers is
\sect{SECTION-II (ii)}
\qn{5} If $(1+ax+bx^2)^4=a_0+a_1x+a_2x^2+\ldots+a_8x^8$ ; where a, b, $a_0$, $a_1$, \ldots\ $a_8\in R$ and $a_0+a_1+a_2\ne0$ and
\[\begin{vmatrix}a_0&a_1&a_2\\a_1&a_2&a_0\\a_2&a_0&a_1\end{vmatrix}=0\]
then value of $\dfrac{5a}{b}$ is equal to
\qn{6} The line $x+2y+c=0$ intersects the circle $x^2+y^2-4=0$ at two distinct points P and Q. Another line $12x-6y=41$ intersects the circle $x^2+y^2-4x-2y+1=0$ at two distinct points R and S. Find the value of c for which the four points P, Q, R and S are concyclic.
\qn{7} The value of $\displaystyle\lim_{n\to\infty}\left(\sum_{r=1}^{n}\left(\dfrac{2r+5}{r^2+r}\right)\left(\dfrac{3}{5}\right)^{r+1}\right)$ is equal to L, then the value of 10L is
\qn{8} If the maximum value of
\[\left|\sqrt{x^4-7x^2-4x+20}-\sqrt{x^4+9x^2+16}\right|\]
is P, find the value of $\left[\dfrac{P}{2}\right]$\\
(where [.] denotes GIF)
\qn{9} Minimum distance between the parabolas $y^2-4x-8y+40=0$ and $x^2-8x-4y+40=0$, is `d' then $\dfrac{5d}{\sqrt2}$ is equal to
\qn{10} If $g:[1,\infty)\to[2,\infty)$ be a differentiable function and satisfying $6\displaystyle\int_1^x g(t)\,dt=3xg(x)-x^3-5\ \forall\ x\ge1$, then value of $\dfrac12 g(5)$ is equal to :
\end{multicols}
\end{document}
